\section{Estructura de la membrana}

Una membrana biológica es una estructura lipídica que establece el límite entre la célula y el espacio extracelular. Está formada por una capa doble de lípidos anfipáticos, que son llamados así por tener una zona polar y una no polar, dirigiendo la polar hacia dentro de la célula y la no polar hacia fuera.

\section{Lípidos de la membrana}

\subsection{Fosfolípidos}

Los fosfolípidos son los principales componentes de la membrana. Están formados por:
\begin{itemize}
    \item Glicerol como base
    \item 2 colas de ácidos grasos (no polares)
    \item Cabeza de grupo fosfato (polar)
\end{itemize}

Su estructura anfipática permite la formación espontánea de bicapas en medio acuoso.

\subsection{Colesterol}

El colesterol es un lípido esteroide que:
\begin{itemize}
    \item Está formado por cuatro anillos de carbono fusionados
    \item Se intercala entre los fosfolípidos
    \item Regula la fluidez de la membrana
    \item Aumenta la estabilidad estructural
\end{itemize}

\section{Incrustaciones}

\subsection{Carbohidratos}
\begin{itemize}
    \item Se encuentran solo en la cara externa de la membrana
    \item Se unen formando glucoproteínas y glucolípidos
    \item Forman el glicocáliz o "capa de azúcar"
    \item Funciones: Reconocimiento celular, adhesión, protección
\end{itemize}

\subsection{Proteínas}
Las proteínas de membrana son incrustadas parcial o totalmente en ambos lados o a través de la membrana. Pueden ser:

\begin{itemize}
    \item \textbf{Integrales}: Atraviesan completamente la membrana
    \item \textbf{Periféricas}: Se asocian débilmente a la superficie
    \item \textbf{Ancladas a lípidos}: Unidas covalentemente a lípidos
\end{itemize}

\subsection{Funciones de las proteínas de membrana}
\begin{itemize}
    \item Transporte de sustancias
    \item Actividad enzimática
    \item Transducción de señales
    \item Adhesión celular
    \item Reconocimiento celular
\end{itemize}